% Options for packages loaded elsewhere
\PassOptionsToPackage{unicode}{hyperref}
\PassOptionsToPackage{hyphens}{url}
\documentclass[
]{article}
\usepackage{xcolor}
\usepackage[margin=1in]{geometry}
\usepackage{amsmath,amssymb}
\setcounter{secnumdepth}{-\maxdimen} % remove section numbering
\usepackage{iftex}
\ifPDFTeX
  \usepackage[T1]{fontenc}
  \usepackage[utf8]{inputenc}
  \usepackage{textcomp} % provide euro and other symbols
\else % if luatex or xetex
  \usepackage{unicode-math} % this also loads fontspec
  \defaultfontfeatures{Scale=MatchLowercase}
  \defaultfontfeatures[\rmfamily]{Ligatures=TeX,Scale=1}
\fi
\usepackage{lmodern}
\ifPDFTeX\else
  % xetex/luatex font selection
\fi
% Use upquote if available, for straight quotes in verbatim environments
\IfFileExists{upquote.sty}{\usepackage{upquote}}{}
\IfFileExists{microtype.sty}{% use microtype if available
  \usepackage[]{microtype}
  \UseMicrotypeSet[protrusion]{basicmath} % disable protrusion for tt fonts
}{}
\makeatletter
\@ifundefined{KOMAClassName}{% if non-KOMA class
  \IfFileExists{parskip.sty}{%
    \usepackage{parskip}
  }{% else
    \setlength{\parindent}{0pt}
    \setlength{\parskip}{6pt plus 2pt minus 1pt}}
}{% if KOMA class
  \KOMAoptions{parskip=half}}
\makeatother
\usepackage{color}
\usepackage{fancyvrb}
\newcommand{\VerbBar}{|}
\newcommand{\VERB}{\Verb[commandchars=\\\{\}]}
\DefineVerbatimEnvironment{Highlighting}{Verbatim}{commandchars=\\\{\}}
% Add ',fontsize=\small' for more characters per line
\usepackage{framed}
\definecolor{shadecolor}{RGB}{248,248,248}
\newenvironment{Shaded}{\begin{snugshade}}{\end{snugshade}}
\newcommand{\AlertTok}[1]{\textcolor[rgb]{0.94,0.16,0.16}{#1}}
\newcommand{\AnnotationTok}[1]{\textcolor[rgb]{0.56,0.35,0.01}{\textbf{\textit{#1}}}}
\newcommand{\AttributeTok}[1]{\textcolor[rgb]{0.13,0.29,0.53}{#1}}
\newcommand{\BaseNTok}[1]{\textcolor[rgb]{0.00,0.00,0.81}{#1}}
\newcommand{\BuiltInTok}[1]{#1}
\newcommand{\CharTok}[1]{\textcolor[rgb]{0.31,0.60,0.02}{#1}}
\newcommand{\CommentTok}[1]{\textcolor[rgb]{0.56,0.35,0.01}{\textit{#1}}}
\newcommand{\CommentVarTok}[1]{\textcolor[rgb]{0.56,0.35,0.01}{\textbf{\textit{#1}}}}
\newcommand{\ConstantTok}[1]{\textcolor[rgb]{0.56,0.35,0.01}{#1}}
\newcommand{\ControlFlowTok}[1]{\textcolor[rgb]{0.13,0.29,0.53}{\textbf{#1}}}
\newcommand{\DataTypeTok}[1]{\textcolor[rgb]{0.13,0.29,0.53}{#1}}
\newcommand{\DecValTok}[1]{\textcolor[rgb]{0.00,0.00,0.81}{#1}}
\newcommand{\DocumentationTok}[1]{\textcolor[rgb]{0.56,0.35,0.01}{\textbf{\textit{#1}}}}
\newcommand{\ErrorTok}[1]{\textcolor[rgb]{0.64,0.00,0.00}{\textbf{#1}}}
\newcommand{\ExtensionTok}[1]{#1}
\newcommand{\FloatTok}[1]{\textcolor[rgb]{0.00,0.00,0.81}{#1}}
\newcommand{\FunctionTok}[1]{\textcolor[rgb]{0.13,0.29,0.53}{\textbf{#1}}}
\newcommand{\ImportTok}[1]{#1}
\newcommand{\InformationTok}[1]{\textcolor[rgb]{0.56,0.35,0.01}{\textbf{\textit{#1}}}}
\newcommand{\KeywordTok}[1]{\textcolor[rgb]{0.13,0.29,0.53}{\textbf{#1}}}
\newcommand{\NormalTok}[1]{#1}
\newcommand{\OperatorTok}[1]{\textcolor[rgb]{0.81,0.36,0.00}{\textbf{#1}}}
\newcommand{\OtherTok}[1]{\textcolor[rgb]{0.56,0.35,0.01}{#1}}
\newcommand{\PreprocessorTok}[1]{\textcolor[rgb]{0.56,0.35,0.01}{\textit{#1}}}
\newcommand{\RegionMarkerTok}[1]{#1}
\newcommand{\SpecialCharTok}[1]{\textcolor[rgb]{0.81,0.36,0.00}{\textbf{#1}}}
\newcommand{\SpecialStringTok}[1]{\textcolor[rgb]{0.31,0.60,0.02}{#1}}
\newcommand{\StringTok}[1]{\textcolor[rgb]{0.31,0.60,0.02}{#1}}
\newcommand{\VariableTok}[1]{\textcolor[rgb]{0.00,0.00,0.00}{#1}}
\newcommand{\VerbatimStringTok}[1]{\textcolor[rgb]{0.31,0.60,0.02}{#1}}
\newcommand{\WarningTok}[1]{\textcolor[rgb]{0.56,0.35,0.01}{\textbf{\textit{#1}}}}
\usepackage{longtable,booktabs,array}
\usepackage{calc} % for calculating minipage widths
% Correct order of tables after \paragraph or \subparagraph
\usepackage{etoolbox}
\makeatletter
\patchcmd\longtable{\par}{\if@noskipsec\mbox{}\fi\par}{}{}
\makeatother
% Allow footnotes in longtable head/foot
\IfFileExists{footnotehyper.sty}{\usepackage{footnotehyper}}{\usepackage{footnote}}
\makesavenoteenv{longtable}
\usepackage{graphicx}
\makeatletter
\newsavebox\pandoc@box
\newcommand*\pandocbounded[1]{% scales image to fit in text height/width
  \sbox\pandoc@box{#1}%
  \Gscale@div\@tempa{\textheight}{\dimexpr\ht\pandoc@box+\dp\pandoc@box\relax}%
  \Gscale@div\@tempb{\linewidth}{\wd\pandoc@box}%
  \ifdim\@tempb\p@<\@tempa\p@\let\@tempa\@tempb\fi% select the smaller of both
  \ifdim\@tempa\p@<\p@\scalebox{\@tempa}{\usebox\pandoc@box}%
  \else\usebox{\pandoc@box}%
  \fi%
}
% Set default figure placement to htbp
\def\fps@figure{htbp}
\makeatother
\setlength{\emergencystretch}{3em} % prevent overfull lines
\providecommand{\tightlist}{%
  \setlength{\itemsep}{0pt}\setlength{\parskip}{0pt}}
\usepackage{bookmark}
\IfFileExists{xurl.sty}{\usepackage{xurl}}{} % add URL line breaks if available
\urlstyle{same}
\hypersetup{
  pdftitle={Lettuce\_growth\_days},
  pdfauthor={Zelimkhan},
  hidelinks,
  pdfcreator={LaTeX via pandoc}}

\title{Lettuce\_growth\_days}
\author{Zelimkhan}
\date{2025-09-22}

\begin{document}
\maketitle

{
\setcounter{tocdepth}{5}
\tableofcontents
}
\section{Install packages and attach
libraries}\label{install-packages-and-attach-libraries}

\begin{Shaded}
\begin{Highlighting}[]
\CommentTok{\# install.packages("psych")}
\CommentTok{\# install.packages("lubridate")}
\CommentTok{\# install.packages("purrr")}
\CommentTok{\# install.packages("GGally")}
\CommentTok{\# install.packages("skimr")}
\CommentTok{\# install.packages("Writexl")}
\FunctionTok{library}\NormalTok{(psych)}
\FunctionTok{library}\NormalTok{(googlesheets4)}
\FunctionTok{library}\NormalTok{(tidyr)}
\FunctionTok{library}\NormalTok{(dplyr)}
\FunctionTok{library}\NormalTok{(janitor)}
\FunctionTok{library}\NormalTok{(ggplot2)}
\FunctionTok{library}\NormalTok{(lubridate)}
\FunctionTok{library}\NormalTok{(purrr)}
\FunctionTok{library}\NormalTok{(GGally)}
\FunctionTok{library}\NormalTok{(readxl)}
\FunctionTok{library}\NormalTok{(skimr)}
\FunctionTok{library}\NormalTok{(writexl)}
\end{Highlighting}
\end{Shaded}

\section{Step 1. Ask}\label{step-1.-ask}

\subsection{1.1 Key tasks and
Deliverable(s)}\label{key-tasks-and-deliverables}

\subsubsection{1.1.1 A clear statement of the business
task}\label{a-clear-statement-of-the-business-task}

Determine how environmental factors (temperature, humidity, and total
dissolved solids) influence lettuce growth and maturity time, in order
to optimize greenhouse conditions and predict growth outcomes for future
growing seasons.

\subsubsection{1.1.2 Key stakeholders}\label{key-stakeholders}

Managers

\subsubsection{1.1.3 Smart questions}\label{smart-questions}

\begin{longtable}[]{@{}
  >{\raggedright\arraybackslash}p{(\linewidth - 4\tabcolsep) * \real{0.0435}}
  >{\raggedright\arraybackslash}p{(\linewidth - 4\tabcolsep) * \real{0.4783}}
  >{\raggedright\arraybackslash}p{(\linewidth - 4\tabcolsep) * \real{0.4783}}@{}}
\toprule\noalign{}
\begin{minipage}[b]{\linewidth}\raggedright
N
\end{minipage} & \begin{minipage}[b]{\linewidth}\raggedright
Type
\end{minipage} & \begin{minipage}[b]{\linewidth}\raggedright
Questions
\end{minipage} \\
\midrule\noalign{}
\endhead
\bottomrule\noalign{}
\endlastfoot
& & \\
1 & Specific & Is there a connection between humidity and
temperature? \\
& & \\
2 & Measurable & How much growth is normal or expected over a one week
period? \\
& & \\
3 & Action-oriented & If total dissolved solids are increased when
growing lettuce, will it shorten the time needed for the plants to reach
maturity? \\
& & \\
4 & Relevant & Based on the dataset, how can predictions about growth be
made given the conditions? \\
& & \\
5 & Time - bound & 1) What is the data range for this data set? 2) Do we
need an understanding of this data for next season? \\
\end{longtable}

\subsubsection{1.1.4 Scope of work (SOW)}\label{scope-of-work-sow}

\begin{longtable}[]{@{}
  >{\raggedright\arraybackslash}p{(\linewidth - 10\tabcolsep) * \real{0.0192}}
  >{\raggedright\arraybackslash}p{(\linewidth - 10\tabcolsep) * \real{0.2115}}
  >{\raggedright\arraybackslash}p{(\linewidth - 10\tabcolsep) * \real{0.0769}}
  >{\raggedright\arraybackslash}p{(\linewidth - 10\tabcolsep) * \real{0.2308}}
  >{\raggedright\arraybackslash}p{(\linewidth - 10\tabcolsep) * \real{0.2308}}
  >{\raggedright\arraybackslash}p{(\linewidth - 10\tabcolsep) * \real{0.2308}}@{}}
\toprule\noalign{}
\begin{minipage}[b]{\linewidth}\raggedright
N
\end{minipage} & \begin{minipage}[b]{\linewidth}\raggedright
Deliverables
\end{minipage} & \begin{minipage}[b]{\linewidth}\raggedright
Timeline
\end{minipage} & \begin{minipage}[b]{\linewidth}\raggedright
Milestones
\end{minipage} & \begin{minipage}[b]{\linewidth}\raggedright
Results
\end{minipage} & \begin{minipage}[b]{\linewidth}\raggedright
What SMART questions will be answered
\end{minipage} \\
\midrule\noalign{}
\endhead
\bottomrule\noalign{}
\endlastfoot
& & & & & \\
1 & Clean dataset & Day 1 & Completion of data download and preparation
& Dataset is clean and usable & - \\
& & & & & \\
2 & Exploratory data analysis (EDA) findings (form: spreadsheet, google
doc summary) & Day 2 & Completion of EDA - understanding of dataset is
formed & EDA highlights room for investigation. Questions were
formulated based on the results of EDA. & \textbf{Specific:} Is the
connection between humidity and temperature \textbf{Time - Bound:} What
is the data range for this data set? \\
& & & & & \\
3 & Engineered features in the form of new columns (examples may include
Avg 3 day growth, Growth since day before, etc.) & Day 3 & Completion of
feature engineering & New columns engineered to provide more accessible
information & \textbf{Measurable:} How much growth is normal or expected
over a one week period? \\
& & & & & \\
4 & Statistical insights & Day 3 & Completion of statistical analysis &
Statistical analysis conducted & \textbf{Action-oriented:} If total
dissolved solids are increased when growing lettuce, will it shorten the
time needed for the plants to reach maturity? \textbf{Relevant:} Based
on the dataset, how can predictions about growth be made given the
conditions? \textbf{Time - bound:} Do we need an understanding of this
data for next season? \\
& & & & & \\
5 & Visual Aids (images, dashboards) & Days 4-5 & Completion of
visualizations & Visualizations created & - \\
& & & & & \\
6 & Final report (spreadsheet and PowerPoint) & Day 6-10 & Final report
submission & Final report completed & - \\
\end{longtable}

\subsection{1.2 Guiding questions}\label{guiding-questions}

\begin{longtable}[]{@{}
  >{\raggedright\arraybackslash}p{(\linewidth - 4\tabcolsep) * \real{0.0435}}
  >{\raggedright\arraybackslash}p{(\linewidth - 4\tabcolsep) * \real{0.4783}}
  >{\raggedright\arraybackslash}p{(\linewidth - 4\tabcolsep) * \real{0.4783}}@{}}
\toprule\noalign{}
\begin{minipage}[b]{\linewidth}\raggedright
N
\end{minipage} & \begin{minipage}[b]{\linewidth}\raggedright
Questions
\end{minipage} & \begin{minipage}[b]{\linewidth}\raggedright
Answers
\end{minipage} \\
\midrule\noalign{}
\endhead
\bottomrule\noalign{}
\endlastfoot
& & \\
1 & What is the problem you are trying to solve? & The challenge lies in
understanding (finding patterns) how environmental conditions influence
lettuce growth time. This knowledge will help inform decisions to
optimize growth cycle. \\
& & \\
2 & How can you insights drive business decisions? & Knowing about
environmental conditions that impact lettuce growth can help managers
decide on construction of an appropriate greenhouse \\
\end{longtable}

\section{Step 2. Prepare}\label{step-2.-prepare}

\subsection{2.1 Key tasks and
deliverable(s)}\label{key-tasks-and-deliverables-1}

\subsubsection{2.1.1 Download data and store it
appropriately}\label{download-data-and-store-it-appropriately}

Data downloaded and stored on my personal computer and Google Drive.
Data available at these links: 1) On the work computer:
C:/Users/GudievZK/Desktop/GitHub/Google Data Analytics Professional
Certificate/1.
Ask\_Questions\_to\_Make\_Data\_Driven\_Decisions/lettuce\_dataset\_SOW.xlsx,
sheet = ``source\_lettuce\_dataset 2) On the laptop:
/Users/zelimkhan/Desktop/Data/GitHub/Google Data Analytics Professional
Certificate/1.
Ask\_Questions\_to\_Make\_Data\_Driven\_Decisions/lettuce\_dataset\_SOW.xlsx
3) On the Google Drive

\subsubsection{2.1.2 Identify how it's
orginized}\label{identify-how-its-orginized}

The data is organized as a table containing data on 70 plants. The rows
contain environmental conditions for each of the 70 plants on each day
throughout the growing season. The columns contain information about the
plant identifier (Plant\_ID) and the following environmental conditions,
which were recorded daily for each of the 70 plants throughout the
growing season: Date; Temperature (?C); Humidity (\%); TDS Value (ppm);
pH Level; Growth Days.

\subsubsection{2.1.3 Sort and filter the
data}\label{sort-and-filter-the-data}

Data filtered and sorted

\subsubsection{2.1.4 Determine the credibility of the
data}\label{determine-the-credibility-of-the-data}

The following methodologies/frameworks were used to determine the
credibility of the data: ROCCC; Data Integrity.

\paragraph{2.1.4.1 ROCCC}\label{roccc}

\begin{longtable}[]{@{}
  >{\raggedright\arraybackslash}p{(\linewidth - 8\tabcolsep) * \real{0.0222}}
  >{\raggedright\arraybackslash}p{(\linewidth - 8\tabcolsep) * \real{0.2444}}
  >{\raggedright\arraybackslash}p{(\linewidth - 8\tabcolsep) * \real{0.2444}}
  >{\raggedright\arraybackslash}p{(\linewidth - 8\tabcolsep) * \real{0.2444}}
  >{\raggedright\arraybackslash}p{(\linewidth - 8\tabcolsep) * \real{0.2444}}@{}}
\toprule\noalign{}
\begin{minipage}[b]{\linewidth}\raggedright
N
\end{minipage} & \begin{minipage}[b]{\linewidth}\raggedright
ROCCC Dimension
\end{minipage} & \begin{minipage}[b]{\linewidth}\raggedright
Description
\end{minipage} & \begin{minipage}[b]{\linewidth}\raggedright
Conclusion
\end{minipage} & \begin{minipage}[b]{\linewidth}\raggedright
Comments
\end{minipage} \\
\midrule\noalign{}
\endhead
\bottomrule\noalign{}
\endlastfoot
1 & Reliable & Source is dependable and consistent. & Watch conclusion
below the table & It depends on ROCCC dimensions listed below \\
2 & Original & We get data from original source (or If we obtained data
from a secondary or tertiary source, we could compare and verify it with
the primary source) & The data is not original & We do not receive data
from the original source, we also cannot identify the original source
and have no way to compare it with the original data. However, since the
data is provided and used for educational purposes, we may use it for
educational purposes \\
3 & Comprehensive & Covers all necessary cases and variables. & It is
not possible to determine the level of comprehensiveness of the data &
It is logical to assume that data is greatly simplified for educational
purposes and does not contain all necessary information needed to solve
a business problem. Since it includes the following key variables it can
be used for our educational purposes: Plant\_ID; Date; Temperature (°C);
Humidity (\%); TDS Value (ppm); pH Level; Growth Days. \\
4 & Current & The data sources aren't out of date or irrelevant & The
data is current & The data can be considered current, even though the
observation was conducted in 2023 and the analysis using this
information is being carried out in 2025. \\
5 & Cited & Data source is cited and vetted & The source of the data was
not cited or verified. & However, since the data is provided and used
for training purposes the data can be used in our work. \\
\end{longtable}

\textbf{Reliable:} The data isn't reliable

\textbf{Conclusion on the credibility of the data in accordance with the
ROCCC framework:} Data isn't credible.

\paragraph{2.1.4.2 Data integrity}\label{data-integrity}

\begin{longtable}[]{@{}
  >{\raggedright\arraybackslash}p{(\linewidth - 8\tabcolsep) * \real{0.0222}}
  >{\raggedright\arraybackslash}p{(\linewidth - 8\tabcolsep) * \real{0.2444}}
  >{\raggedright\arraybackslash}p{(\linewidth - 8\tabcolsep) * \real{0.2444}}
  >{\raggedright\arraybackslash}p{(\linewidth - 8\tabcolsep) * \real{0.2444}}
  >{\raggedright\arraybackslash}p{(\linewidth - 8\tabcolsep) * \real{0.2444}}@{}}
\toprule\noalign{}
\begin{minipage}[b]{\linewidth}\raggedright
N
\end{minipage} & \begin{minipage}[b]{\linewidth}\raggedright
Data integrity Dimension
\end{minipage} & \begin{minipage}[b]{\linewidth}\raggedright
Description
\end{minipage} & \begin{minipage}[b]{\linewidth}\raggedright
Conclusion
\end{minipage} & \begin{minipage}[b]{\linewidth}\raggedright
Comments
\end{minipage} \\
\midrule\noalign{}
\endhead
\bottomrule\noalign{}
\endlastfoot
1 & Accuracy & The degree to which the data conforms to the actual
entity being measured or described & It is not possible to determine the
data's accuracy. & Since the data does not contain unrealistic values
(such as days with negative growth, numerical values falling within an
unexpected data range, or other impossible measurements), the data can
be considered accurate. \\
2 & Completeness & Data completeness is a dimension of data quality that
measures whether all necessary or required data is present and available
for a dataset to fulfill its intended purpose. & It is not possible to
determine the level of Completeness of the data & It is logical to
assume that data is greatly simplified for educational purposes and does
not contain all necessary information needed to solve a business
problem. Since it includes the following key variables and doesn't have
missing values it can be used for our educational purposes: Plant\_ID;
Date; Temperature (°C); Humidity (\%); TDS Value (ppm); pH Level; Growth
Days. \\
3 & Consistency & Data follows the same format and rules. & Data is
consistency & The Date column has character type. \\
4 & Trustworthiness & Data trustworthiness refers to the degree to which
information can be relied upon as accurate, credible, and secure,
allowing users to confidently make decisions and conduct analyses based
on it. & The data isn't trustworthiness & The data isn't trustworthy due
to low scores for data integrity dimensions listed above. \\
\end{longtable}

\textbf{Trustworthiness:} The data isn't trustworthiness.

\textbf{Conclusion on credibility of the data in accordance with the
Data integrity conception:} The data isn't credible.

\paragraph{2.1.4.3 Bias}\label{bias}

The dataset appears to come from a controlled educational or research
environment, so it may not present all real world farming conditions or
varieties of lettuce. However, since it is used for analytical practice
rather than policy or production decisions, the risk of bias affecting
conclusions is minimal.

\textbf{Conclusion on the admissibility of using data:} The data are
acceptable for use despite negative findings regarding the credibility
of the data in accordance with both the ROCCC framework and the concept
of data integrity, as well as the inability to demonstrate the absence
of bias, as they are used for educational purposes.

\subsubsection{2.1.5 A description of all data sourced
used}\label{a-description-of-all-data-sourced-used}

For this analytical task, the ``lettuce growth'' dataset from Kaggle was
used.

\subsection{2.2 Guiding questions}\label{guiding-questions-1}

\begin{longtable}[]{@{}
  >{\raggedright\arraybackslash}p{(\linewidth - 4\tabcolsep) * \real{0.0435}}
  >{\raggedright\arraybackslash}p{(\linewidth - 4\tabcolsep) * \real{0.4783}}
  >{\raggedright\arraybackslash}p{(\linewidth - 4\tabcolsep) * \real{0.4783}}@{}}
\toprule\noalign{}
\begin{minipage}[b]{\linewidth}\raggedright
N
\end{minipage} & \begin{minipage}[b]{\linewidth}\raggedright
Questions
\end{minipage} & \begin{minipage}[b]{\linewidth}\raggedright
Answers
\end{minipage} \\
\midrule\noalign{}
\endhead
\bottomrule\noalign{}
\endlastfoot
& & \\
1 & Where is your data located? & The data is publicly available in
Kaggle website at this
\href{https://www.kaggle.com/datasets/jurijsruko/lettuce}{link} \\
& & \\
2 & How is the data organized? & The data is organized as a table
containing data on 70 plants in long format. The rows contain
environmental conditions for each of the 70 plants on each day
throughout the growing season. The columns contain information about the
plant identifier (Plant\_ID) and the following environmental conditions,
which were recorded daily for each of the 70 plants throughout the
growing season: Date; Temperature (?C); Humidity (\%); TDS Value (ppm);
pH Level; Growth Days. \\
& & \\
3 & Are there issues with bias or credibility in this data? Does your
data ROCCC? & See section ``2.1.4 Determine the credibility of the
data'' \\
& & \\
4 & How are you addressing licensing, privacy, security, and
accessibility? & 1. \textbf{Licensing} The lettuce dataset is published
under the Apache 2.0.It means: 1) You are allowed to freely use, modify,
and share the dataset, including for commercial and non-commercial
purposes. 2) You must provide proper attribution to the original authors
when you use the data. 3) If you modify the dataset and distribute your
modified version, you must indicate what changes you made. 4) There is
no warranty---the authors are not responsible for any issues from using
the data. 4) The dataset can be combined with other work, even
proprietary. However, on the Kaggle page for this dataset, there is an
additional statement: ``Do not use it for commercial purposes.'' This
suggests the author is placing an extra restriction that is stricter
than the typical Apache 2.0 terms. In this case, you should follow the
conditions stated on the dataset's page. For academic or personal
projects, with attribution, you are free to use the data. 2.
\textbf{Privacy} It does not contain personally identifiable information
(PII), so privacy is not a concern. 3. \textbf{Security} Security risks
are minimal, as the dataset stored locally on my Google Drive. 4.
\textbf{Accessibility} Accessibility is ensured by keeping the data in
standard formats (csv, xlsx, R data frames), which can be opened in
widely available tools such as Excel, Rstudio or Google Sheets. \\
& & \\
5 & How did you verify the data's integrity? & Data integrity involves
the accuracy, completeness, consistency, and trustworthiness of data
throughout its lifecycle. 1. \textbf{Accuracy:} I inspected the dataset
for unrealistic values (e.g., negative growth days, numerical values
falling within an unexpected data range or impossible measurements)
using: Sorting and filtering in spreadsheet; R functions such as
summary() and is.na(). \emph{Definition Data accuracy refers to the
degree to which data correctly represents the true, real-world
information or events it describes, meaning it is correct, precise, and
free from errors or distortions. Accurate data serves as a trustworthy
foundation for effective decision-making, reliable analysis, and
efficient operations, making it a critical component of overall data
quality.} 2. \textbf{Completeness:} I checked for missing or empty cells
using: Sorting and filtering in spreadsheet; R functions such as
summary() and is.na(). \emph{Definition Data completeness is a dimension
of data quality that measures whether all necessary or required data is
present and available for a dataset to fulfill its intended purpose. It
ensures there are no missing values or gaps in critical information,
preventing skewed analyses and enabling accurate decision-making}. 3.
\textbf{Consistency:} I reviewed column names, formats, and data types
to ensure that values were stored uniformly (e.g., all numeric columns
have numeric data types). \emph{Definition Data consistency - is the
quality of data follows the same format and rules remaining accurate,
uniform, and coherent across all systems, databases, and applications
within an organization.}. 4. \textbf{Trustworthiness:} Users can depend
on the data to be correct and dependable. Since the dataset is a clean
educational sample, I verified its source and confirmed it had not been
altered during download or import (no unexpected changes in row count or
structure). \emph{Definition Data trustworthiness refers to the degree
to which information can be relied upon as accurate, credible, and
secure, allowing users to confidently make decisions and conduct
analyses based on it.} Data?s integrity will be checked in more detail
(accuracy, completeness, consistency, trustworthiness) in the next step
of analyses (step 3 ?Process?). \\
& & \\
6 & How does it help you answer you questions? & The dataset directly
contains the variables needed (growth days and environmental factors),
so it enables me to explore how growth varies by environmental
conditions). It provides the empirical basis for descriptive statistics
and visualizations that can answer the practice assignment's research
questions. \\
& & \\
7 & Are there any problems with the data? & The Date column has
character type instead date. \\
\end{longtable}

\section{Step 3. Process (EDA)}\label{step-3.-process-eda}

\subsection{Stage 1 of EDA: Data Discovery \&
Profiling}\label{stage-1-of-eda-data-discovery-profiling}

\paragraph{Load data}\label{load-data}

\begin{Shaded}
\begin{Highlighting}[]
\NormalTok{df\_0 }\OtherTok{\textless{}{-}} \FunctionTok{read\_sheet}\NormalTok{(}\StringTok{"https://docs.google.com/spreadsheets/d/1esw9BCXZDO1h4hwuaKowia5eVmflEKTS6NiouamSYA0/edit?gid=723124268\#gid=723124268"}\NormalTok{, }\AttributeTok{sheet =} \StringTok{"source\_lettuce\_dataset"}\NormalTok{)}

\CommentTok{\# df\_0 \textless{}{-} read\_excel("C:/Users/GudievZK/Desktop/GitHub/Google Data Analytics Professional Certificate/1. Ask\_Questions\_to\_Make\_Data\_Driven\_Decisions/lettuce\_dataset\_SOW.xlsx", sheet = "source\_lettuce\_dataset")}
\end{Highlighting}
\end{Shaded}

\paragraph{Table overview}\label{table-overview}

\begin{Shaded}
\begin{Highlighting}[]
\NormalTok{df\_0}
\end{Highlighting}
\end{Shaded}

\begin{verbatim}
## # A tibble: 3,169 x 7
##    Plant_ID Date  `Temperature (°C)` `Humidity (%)` `TDS Value (ppm)` `pH Level`
##       <dbl> <chr>              <dbl>          <dbl>             <dbl>      <dbl>
##  1        1 9/12~               31.2             62               576        6.7
##  2        2 9/12~               31.2             66               657        6.8
##  3        3 9/12~               31.2             69               403        6.6
##  4        4 9/12~               31.2             72               729        6.6
##  5        5 9/12~               31.2             72               738        6.2
##  6        6 9/12~               31.2             57               793        6.4
##  7        7 9/12~               31.2             52               468        6.7
##  8        8 9/12~               31.2             51               768        6.1
##  9        9 9/12~               31.2             80               793        6.6
## 10       10 9/12~               31.2             58               794        6.6
## # i 3,159 more rows
## # i 1 more variable: `Growth Days` <dbl>
\end{verbatim}

\paragraph{Column names}\label{column-names}

\begin{Shaded}
\begin{Highlighting}[]
\FunctionTok{colnames}\NormalTok{(df\_0)}
\end{Highlighting}
\end{Shaded}

\begin{verbatim}
## [1] "Plant_ID"         "Date"             "Temperature (°C)" "Humidity (%)"    
## [5] "TDS Value (ppm)"  "pH Level"         "Growth Days"
\end{verbatim}

\paragraph{Number of rows and columns}\label{number-of-rows-and-columns}

\begin{Shaded}
\begin{Highlighting}[]
\FunctionTok{dim}\NormalTok{(df\_0)}
\end{Highlighting}
\end{Shaded}

\begin{verbatim}
## [1] 3169    7
\end{verbatim}

\paragraph{Data types}\label{data-types}

\begin{Shaded}
\begin{Highlighting}[]
\FunctionTok{glimpse}\NormalTok{(df\_0)}
\end{Highlighting}
\end{Shaded}

\begin{verbatim}
## Rows: 3,169
## Columns: 7
## $ Plant_ID           <dbl> 1, 2, 3, 4, 5, 6, 7, 8, 9, 10, 11, 12, 13, 14, 15, ~
## $ Date               <chr> "9/12/2023", "9/12/2023", "9/12/2023", "9/12/2023",~
## $ `Temperature (°C)` <dbl> 31.2, 31.2, 31.2, 31.2, 31.2, 31.2, 31.2, 31.2, 31.~
## $ `Humidity (%)`     <dbl> 62, 66, 69, 72, 72, 57, 52, 51, 80, 58, 66, 64, 51,~
## $ `TDS Value (ppm)`  <dbl> 576, 657, 403, 729, 738, 793, 468, 768, 793, 794, 4~
## $ `pH Level`         <dbl> 6.7, 6.8, 6.6, 6.6, 6.2, 6.4, 6.7, 6.1, 6.6, 6.6, 6~
## $ `Growth Days`      <dbl> 41, 41, 41, 41, 41, 41, 41, 41, 41, 41, 41, 41, 41,~
\end{verbatim}

\paragraph{Data classes}\label{data-classes}

\begin{Shaded}
\begin{Highlighting}[]
\FunctionTok{sapply}\NormalTok{(df\_0, class)}
\end{Highlighting}
\end{Shaded}

\begin{verbatim}
##         Plant_ID             Date Temperature (°C)     Humidity (%) 
##        "numeric"      "character"        "numeric"        "numeric" 
##  TDS Value (ppm)         pH Level      Growth Days 
##        "numeric"        "numeric"        "numeric"
\end{verbatim}

\paragraph{Checking completeness}\label{checking-completeness}

Number of missing values in each column

\begin{Shaded}
\begin{Highlighting}[]
\FunctionTok{colSums}\NormalTok{(}\FunctionTok{is.na}\NormalTok{(df\_0))}
\end{Highlighting}
\end{Shaded}

\begin{verbatim}
##         Plant_ID             Date Temperature (°C)     Humidity (%) 
##                0                0                0                0 
##  TDS Value (ppm)         pH Level      Growth Days 
##                0                0                0
\end{verbatim}

\paragraph{Checking duplicates}\label{checking-duplicates}

Number of duplicate values in each column

\begin{Shaded}
\begin{Highlighting}[]
\FunctionTok{sapply}\NormalTok{(df\_0, }\ControlFlowTok{function}\NormalTok{(x) }\FunctionTok{sum}\NormalTok{(}\FunctionTok{duplicated}\NormalTok{(x)))}
\end{Highlighting}
\end{Shaded}

\begin{verbatim}
##         Plant_ID             Date Temperature (°C)     Humidity (%) 
##             3099             3121             3071             3138 
##  TDS Value (ppm)         pH Level      Growth Days 
##             2768             3160             3121
\end{verbatim}

\paragraph{Checking uniqueness}\label{checking-uniqueness}

Number of unique rows

\begin{Shaded}
\begin{Highlighting}[]
\FunctionTok{distinct}\NormalTok{(df\_0) }\SpecialCharTok{\%\textgreater{}\%} \FunctionTok{dim}\NormalTok{()}
\end{Highlighting}
\end{Shaded}

\begin{verbatim}
## [1] 3169    7
\end{verbatim}

Number of unique values in each column

\begin{Shaded}
\begin{Highlighting}[]
\FunctionTok{sapply}\NormalTok{(df\_0, }\ControlFlowTok{function}\NormalTok{(x) }\FunctionTok{sum}\NormalTok{(}\FunctionTok{n\_distinct}\NormalTok{(x)))}
\end{Highlighting}
\end{Shaded}

\begin{verbatim}
##         Plant_ID             Date Temperature (°C)     Humidity (%) 
##               70               48               98               31 
##  TDS Value (ppm)         pH Level      Growth Days 
##              401                9               48
\end{verbatim}

\paragraph{Transorm data}\label{transorm-data}

Convert Date column from character into date

\begin{Shaded}
\begin{Highlighting}[]
\NormalTok{df\_0}\SpecialCharTok{$}\NormalTok{Date }\OtherTok{\textless{}{-}} \FunctionTok{mdy}\NormalTok{(df\_0}\SpecialCharTok{$}\NormalTok{Date)}
\end{Highlighting}
\end{Shaded}

Convert Plant\_ID column from numeric into factor

\begin{Shaded}
\begin{Highlighting}[]
\NormalTok{df\_0}\SpecialCharTok{$}\NormalTok{Plant\_ID }\OtherTok{\textless{}{-}} \FunctionTok{as.factor}\NormalTok{(df\_0}\SpecialCharTok{$}\NormalTok{Plant\_ID)}
\end{Highlighting}
\end{Shaded}

Clean names

\begin{Shaded}
\begin{Highlighting}[]
\NormalTok{df }\OtherTok{\textless{}{-}} \FunctionTok{clean\_names}\NormalTok{(df\_0)}
\end{Highlighting}
\end{Shaded}

\paragraph{Checking data after
transormation}\label{checking-data-after-transormation}

Table overview

\begin{Shaded}
\begin{Highlighting}[]
\NormalTok{df }\SpecialCharTok{\%\textgreater{}\%} \FunctionTok{arrange}\NormalTok{(date)}
\end{Highlighting}
\end{Shaded}

\begin{verbatim}
## # A tibble: 3,169 x 7
##    plant_id date       temperature_c humidity_percent tds_value_ppm p_h_level
##    <fct>    <date>             <dbl>            <dbl>         <dbl>     <dbl>
##  1 1        2023-08-03          33.4               53           582       6.4
##  2 2        2023-08-03          33.4               74           586       6.1
##  3 3        2023-08-03          33.4               51           667       6.2
##  4 4        2023-08-03          33.4               62           547       6.7
##  5 5        2023-08-03          33.4               67           651       6.3
##  6 6        2023-08-03          33.4               66           795       6.1
##  7 7        2023-08-03          33.4               74           459       6.5
##  8 8        2023-08-03          33.4               76           650       6  
##  9 9        2023-08-03          33.4               72           693       6.4
## 10 10       2023-08-03          33.4               66           503       6.5
## # i 3,159 more rows
## # i 1 more variable: growth_days <dbl>
\end{verbatim}

Checking completeness after transormation

\begin{Shaded}
\begin{Highlighting}[]
\FunctionTok{colSums}\NormalTok{(}\FunctionTok{is.na}\NormalTok{(df))}
\end{Highlighting}
\end{Shaded}

\begin{verbatim}
##         plant_id             date    temperature_c humidity_percent 
##                0                0                0                0 
##    tds_value_ppm        p_h_level      growth_days 
##                0                0                0
\end{verbatim}

Checking data stricture

\begin{Shaded}
\begin{Highlighting}[]
\FunctionTok{str}\NormalTok{(df)}
\end{Highlighting}
\end{Shaded}

\begin{verbatim}
## tibble [3,169 x 7] (S3: tbl_df/tbl/data.frame)
##  $ plant_id        : Factor w/ 70 levels "1","2","3","4",..: 1 2 3 4 5 6 7 8 9 10 ...
##  $ date            : Date[1:3169], format: "2023-09-12" "2023-09-12" ...
##  $ temperature_c   : num [1:3169] 31.2 31.2 31.2 31.2 31.2 31.2 31.2 31.2 31.2 31.2 ...
##  $ humidity_percent: num [1:3169] 62 66 69 72 72 57 52 51 80 58 ...
##  $ tds_value_ppm   : num [1:3169] 576 657 403 729 738 793 468 768 793 794 ...
##  $ p_h_level       : num [1:3169] 6.7 6.8 6.6 6.6 6.2 6.4 6.7 6.1 6.6 6.6 ...
##  $ growth_days     : num [1:3169] 41 41 41 41 41 41 41 41 41 41 ...
\end{verbatim}

\paragraph{Interpretation of the 1st stage EDA: Data Discovery \&
Profiling}\label{interpretation-of-the-1st-stage-eda-data-discovery-profiling}

\textbf{Purpose and Approach}

This stage focused on initial understanding and validation of the
lettuce dataset to answer two main questions:

What does the data contain?

Is the data clean and ready for analysis?

\textbf{📄 Findings}

Dataset structure: The table contained 3,169 rows and 7 columns,
representing multiple measurements for 70 individual plants across 48
observation days. Columns included: plant\_id, date, temperature\_c,
humidity\_percent, tds\_value\_ppm, p\_h\_level and growth\_days

\textbf{Uniqueness and duplication checks:}

No duplicated rows were found (distinct(df\_0) returned 3,169 unique
rows).

Each variable had a meaningful count of unique values (e.g., 70 plant
IDs, 98 temperature readings, 9 pH levels), confirming valid variability
across measurements.

\textbf{Data type consistency:}

The ``date'' column was initially stored as a character type and was
converted to a proper Date format using mdy().

The ``plant\_id'' column was converted from numeric to factor, which
correctly represents categorical identifiers.

Column names were standardized with clean\_names() for consistent
syntax.

\textbf{Quality assessment:}

No missing values were reported.

The only detected issue before transformation was inconsistent date
formatting.

After cleaning, all variables were in the correct type and ready for
deeper analysis.

\textbf{🧠 Concusion}

The dataset was confirmed to be accurate, consistent, and complete.
Minor formatting issues (dates, variable types) were corrected.

\section{Step 4. Analyze}\label{step-4.-analyze}

\subsection{Stage 2 of EDA: Univariate analysis (``Understanding
individual variables''
phase)}\label{stage-2-of-eda-univariate-analysis-understanding-individual-variables-phase}

\paragraph{Summary statistics}\label{summary-statistics}

Summary statistics using skim function

\begin{Shaded}
\begin{Highlighting}[]
\FunctionTok{skim}\NormalTok{(df) }\SpecialCharTok{\%\textgreater{}\%} 
  \FunctionTok{mutate}\NormalTok{(}\FunctionTok{across}\NormalTok{(}\FunctionTok{where}\NormalTok{(is.numeric), }\SpecialCharTok{\textasciitilde{}} \FunctionTok{round}\NormalTok{(.x, }\DecValTok{2}\NormalTok{)))}
\end{Highlighting}
\end{Shaded}

\begin{longtable}[]{@{}ll@{}}
\caption{Data summary}\tabularnewline
\toprule\noalign{}
\endfirsthead
\endhead
\bottomrule\noalign{}
\endlastfoot
Name & df \\
Number of rows & 3169 \\
Number of columns & 7 \\
\_\_\_\_\_\_\_\_\_\_\_\_\_\_\_\_\_\_\_\_\_\_\_ & \\
Column type frequency: & \\
Date & 1 \\
factor & 1 \\
numeric & 5 \\
\_\_\_\_\_\_\_\_\_\_\_\_\_\_\_\_\_\_\_\_\_\_\_\_ & \\
Group variables & None \\
\end{longtable}

\textbf{Variable type: Date}

\begin{longtable}[]{@{}
  >{\raggedright\arraybackslash}p{(\linewidth - 12\tabcolsep) * \real{0.1750}}
  >{\raggedleft\arraybackslash}p{(\linewidth - 12\tabcolsep) * \real{0.1250}}
  >{\raggedleft\arraybackslash}p{(\linewidth - 12\tabcolsep) * \real{0.1750}}
  >{\raggedright\arraybackslash}p{(\linewidth - 12\tabcolsep) * \real{0.1375}}
  >{\raggedright\arraybackslash}p{(\linewidth - 12\tabcolsep) * \real{0.1375}}
  >{\raggedright\arraybackslash}p{(\linewidth - 12\tabcolsep) * \real{0.1375}}
  >{\raggedleft\arraybackslash}p{(\linewidth - 12\tabcolsep) * \real{0.1125}}@{}}
\toprule\noalign{}
\begin{minipage}[b]{\linewidth}\raggedright
skim\_variable
\end{minipage} & \begin{minipage}[b]{\linewidth}\raggedleft
n\_missing
\end{minipage} & \begin{minipage}[b]{\linewidth}\raggedleft
complete\_rate
\end{minipage} & \begin{minipage}[b]{\linewidth}\raggedright
min
\end{minipage} & \begin{minipage}[b]{\linewidth}\raggedright
max
\end{minipage} & \begin{minipage}[b]{\linewidth}\raggedright
median
\end{minipage} & \begin{minipage}[b]{\linewidth}\raggedleft
n\_unique
\end{minipage} \\
\midrule\noalign{}
\endhead
\bottomrule\noalign{}
\endlastfoot
date & 0 & 1 & 2023-08-03 & 2023-09-19 & 2023-08-25 & 48 \\
\end{longtable}

\textbf{Variable type: factor}

\begin{longtable}[]{@{}
  >{\raggedright\arraybackslash}p{(\linewidth - 10\tabcolsep) * \real{0.1687}}
  >{\raggedleft\arraybackslash}p{(\linewidth - 10\tabcolsep) * \real{0.1205}}
  >{\raggedleft\arraybackslash}p{(\linewidth - 10\tabcolsep) * \real{0.1687}}
  >{\raggedright\arraybackslash}p{(\linewidth - 10\tabcolsep) * \real{0.0964}}
  >{\raggedleft\arraybackslash}p{(\linewidth - 10\tabcolsep) * \real{0.1084}}
  >{\raggedright\arraybackslash}p{(\linewidth - 10\tabcolsep) * \real{0.3373}}@{}}
\toprule\noalign{}
\begin{minipage}[b]{\linewidth}\raggedright
skim\_variable
\end{minipage} & \begin{minipage}[b]{\linewidth}\raggedleft
n\_missing
\end{minipage} & \begin{minipage}[b]{\linewidth}\raggedleft
complete\_rate
\end{minipage} & \begin{minipage}[b]{\linewidth}\raggedright
ordered
\end{minipage} & \begin{minipage}[b]{\linewidth}\raggedleft
n\_unique
\end{minipage} & \begin{minipage}[b]{\linewidth}\raggedright
top\_counts
\end{minipage} \\
\midrule\noalign{}
\endhead
\bottomrule\noalign{}
\endlastfoot
plant\_id & 0 & 1 & FALSE & 70 & 4: 48, 55: 48, 3: 47, 7: 47 \\
\end{longtable}

\textbf{Variable type: numeric}

\begin{longtable}[]{@{}
  >{\raggedright\arraybackslash}p{(\linewidth - 20\tabcolsep) * \real{0.1910}}
  >{\raggedleft\arraybackslash}p{(\linewidth - 20\tabcolsep) * \real{0.1124}}
  >{\raggedleft\arraybackslash}p{(\linewidth - 20\tabcolsep) * \real{0.1573}}
  >{\raggedleft\arraybackslash}p{(\linewidth - 20\tabcolsep) * \real{0.0787}}
  >{\raggedleft\arraybackslash}p{(\linewidth - 20\tabcolsep) * \real{0.0787}}
  >{\raggedleft\arraybackslash}p{(\linewidth - 20\tabcolsep) * \real{0.0449}}
  >{\raggedleft\arraybackslash}p{(\linewidth - 20\tabcolsep) * \real{0.0674}}
  >{\raggedleft\arraybackslash}p{(\linewidth - 20\tabcolsep) * \real{0.0674}}
  >{\raggedleft\arraybackslash}p{(\linewidth - 20\tabcolsep) * \real{0.0674}}
  >{\raggedleft\arraybackslash}p{(\linewidth - 20\tabcolsep) * \real{0.0674}}
  >{\raggedright\arraybackslash}p{(\linewidth - 20\tabcolsep) * \real{0.0674}}@{}}
\toprule\noalign{}
\begin{minipage}[b]{\linewidth}\raggedright
skim\_variable
\end{minipage} & \begin{minipage}[b]{\linewidth}\raggedleft
n\_missing
\end{minipage} & \begin{minipage}[b]{\linewidth}\raggedleft
complete\_rate
\end{minipage} & \begin{minipage}[b]{\linewidth}\raggedleft
mean
\end{minipage} & \begin{minipage}[b]{\linewidth}\raggedleft
sd
\end{minipage} & \begin{minipage}[b]{\linewidth}\raggedleft
p0
\end{minipage} & \begin{minipage}[b]{\linewidth}\raggedleft
p25
\end{minipage} & \begin{minipage}[b]{\linewidth}\raggedleft
p50
\end{minipage} & \begin{minipage}[b]{\linewidth}\raggedleft
p75
\end{minipage} & \begin{minipage}[b]{\linewidth}\raggedleft
p100
\end{minipage} & \begin{minipage}[b]{\linewidth}\raggedright
hist
\end{minipage} \\
\midrule\noalign{}
\endhead
\bottomrule\noalign{}
\endlastfoot
temperature\_c & 0 & 1 & 28.14 & 4.67 & 18 & 23.6 & 30.2 & 31.5 & 33.5 &
▂▂▁▅▇ \\
humidity\_percent & 0 & 1 & 64.87 & 8.99 & 50 & 57.0 & 65.0 & 73.0 &
80.0 & ▇▆▆▇▆ \\
tds\_value\_ppm & 0 & 1 & 598.05 & 115.71 & 400 & 498.0 & 593.0 & 699.0
& 800.0 & ▇▇▇▇▇ \\
p\_h\_level & 0 & 1 & 6.40 & 0.23 & 6 & 6.2 & 6.4 & 6.6 & 6.8 & ▆▇▃▇▆ \\
growth\_days & 0 & 1 & 23.14 & 13.08 & 1 & 12.0 & 23.0 & 34.0 & 48.0 &
▇▇▇▇▆ \\
\end{longtable}

Summary statistics using summary function

\begin{Shaded}
\begin{Highlighting}[]
\FunctionTok{summary}\NormalTok{(df)}
\end{Highlighting}
\end{Shaded}

\begin{verbatim}
##     plant_id         date            temperature_c   humidity_percent
##  4      :  48   Min.   :2023-08-03   Min.   :18.00   Min.   :50.00   
##  55     :  48   1st Qu.:2023-08-14   1st Qu.:23.60   1st Qu.:57.00   
##  3      :  47   Median :2023-08-25   Median :30.20   Median :65.00   
##  7      :  47   Mean   :2023-08-25   Mean   :28.14   Mean   :64.87   
##  11     :  47   3rd Qu.:2023-09-05   3rd Qu.:31.50   3rd Qu.:73.00   
##  48     :  47   Max.   :2023-09-19   Max.   :33.50   Max.   :80.00   
##  (Other):2885                                                        
##  tds_value_ppm   p_h_level      growth_days   
##  Min.   :400   Min.   :6.000   Min.   : 1.00  
##  1st Qu.:498   1st Qu.:6.200   1st Qu.:12.00  
##  Median :593   Median :6.400   Median :23.00  
##  Mean   :598   Mean   :6.399   Mean   :23.14  
##  3rd Qu.:699   3rd Qu.:6.600   3rd Qu.:34.00  
##  Max.   :800   Max.   :6.800   Max.   :48.00  
## 
\end{verbatim}

\paragraph{Examine distribution and detecting
outliers}\label{examine-distribution-and-detecting-outliers}

Examine distribution plotting histograms

\begin{Shaded}
\begin{Highlighting}[]
\NormalTok{df }\SpecialCharTok{\%\textgreater{}\%} \FunctionTok{select}\NormalTok{(}\FunctionTok{where}\NormalTok{(is.numeric)) }\SpecialCharTok{\%\textgreater{}\%} 
  \FunctionTok{pivot\_longer}\NormalTok{(}\FunctionTok{everything}\NormalTok{(),}
               \AttributeTok{names\_to =} \StringTok{"variable"}\NormalTok{,}
               \AttributeTok{values\_to =} \StringTok{"value"}\NormalTok{) }\SpecialCharTok{\%\textgreater{}\%} 
  \FunctionTok{ggplot}\NormalTok{(}\FunctionTok{aes}\NormalTok{(}\AttributeTok{x =}\NormalTok{ value, }\AttributeTok{fill =}\NormalTok{ variable)) }\SpecialCharTok{+}
  \FunctionTok{geom\_histogram}\NormalTok{(}\AttributeTok{bins =} \DecValTok{30}\NormalTok{, }\AttributeTok{alpha =} \FloatTok{0.6}\NormalTok{, }\AttributeTok{color =} \StringTok{"black"}\NormalTok{, }\AttributeTok{na.rm =} \ConstantTok{TRUE}\NormalTok{) }\SpecialCharTok{+}
  \FunctionTok{facet\_wrap}\NormalTok{(}\SpecialCharTok{\textasciitilde{}}\NormalTok{ variable, }\AttributeTok{scales =} \StringTok{"free"}\NormalTok{, }\AttributeTok{ncol =} \DecValTok{3}\NormalTok{) }\SpecialCharTok{+}
\CommentTok{\# theme\_minimal() + }
  \FunctionTok{guides}\NormalTok{(}\AttributeTok{fill =} \StringTok{"none"}\NormalTok{)}
\end{Highlighting}
\end{Shaded}

\pandocbounded{\includegraphics[keepaspectratio]{Lettuce_growth_days_v2.1--Step-2-_files/figure-latex/unnamed-chunk-19-1.pdf}}

Detect outliers plotting boxplots

\begin{Shaded}
\begin{Highlighting}[]
\NormalTok{df }\SpecialCharTok{\%\textgreater{}\%} \FunctionTok{select}\NormalTok{(}\FunctionTok{where}\NormalTok{(is.numeric)) }\SpecialCharTok{\%\textgreater{}\%} 
  \FunctionTok{pivot\_longer}\NormalTok{(}\FunctionTok{everything}\NormalTok{(),}
               \AttributeTok{names\_to =} \StringTok{"variable"}\NormalTok{,}
               \AttributeTok{values\_to =} \StringTok{"value"}\NormalTok{) }\SpecialCharTok{\%\textgreater{}\%} 
  \FunctionTok{ggplot}\NormalTok{(}\FunctionTok{aes}\NormalTok{(}\AttributeTok{x =}\NormalTok{ variable, }\AttributeTok{y =}\NormalTok{ value)) }\SpecialCharTok{+}
  \FunctionTok{geom\_boxplot}\NormalTok{(}\AttributeTok{outliers.alpha =} \FloatTok{0.35}\NormalTok{) }\SpecialCharTok{+}
  \FunctionTok{labs}\NormalTok{(}\AttributeTok{x =} \ConstantTok{NULL}\NormalTok{, }\AttributeTok{y =} \ConstantTok{NULL}\NormalTok{, }\AttributeTok{title =} \StringTok{"Boxplots of numeric variables"}\NormalTok{) }\SpecialCharTok{+}
  \FunctionTok{facet\_wrap}\NormalTok{(}\SpecialCharTok{\textasciitilde{}}\NormalTok{ variable, }\AttributeTok{scales =} \StringTok{"free"}\NormalTok{, }\AttributeTok{ncol =} \DecValTok{3}\NormalTok{) }\SpecialCharTok{+}
  \CommentTok{\# theme\_minimal() +}
  \FunctionTok{guides}\NormalTok{(}\AttributeTok{fill =} \StringTok{"none"}\NormalTok{)}
\end{Highlighting}
\end{Shaded}

\begin{verbatim}
## Warning in geom_boxplot(outliers.alpha = 0.35): Ignoring unknown parameters:
## `outliers.alpha`
\end{verbatim}

\pandocbounded{\includegraphics[keepaspectratio]{Lettuce_growth_days_v2.1--Step-2-_files/figure-latex/unnamed-chunk-20-1.pdf}}

\paragraph{Interpretation of the 2d stage EDA: Univariate
analysis}\label{interpretation-of-the-2d-stage-eda-univariate-analysis}

\textbf{🎯 Purpose}

This step examined each variable independently to describe its
distribution, variation, and potential outliers.

\textbf{📈 Summary Statistics}

Rows: 3,169

Columns: 7

Variable types: 1 Date (date), 1 Factor (plant\_id), 5 Numeric
(temperature\_c, humidity\_percent, tds\_value\_ppm, p\_h\_level,
growth\_days).

Missing data: 0 across all columns.

📅 Date Variable

Range: 2023-08-03 → 2023-09-19 (48 unique days). → Represents about 7
weeks of continuous monitoring.

🪴 Plant ID

70 unique plants, each measured several times (roughly 45 observations
per plant). → Confirms the dataset's panel/time-series nature.

\textbf{🔢 Numeric Variables}

\begin{longtable}[]{@{}
  >{\raggedright\arraybackslash}p{(\linewidth - 4\tabcolsep) * \real{0.0435}}
  >{\raggedright\arraybackslash}p{(\linewidth - 4\tabcolsep) * \real{0.4783}}
  >{\raggedright\arraybackslash}p{(\linewidth - 4\tabcolsep) * \real{0.4783}}@{}}
\toprule\noalign{}
\begin{minipage}[b]{\linewidth}\raggedright
N
\end{minipage} & \begin{minipage}[b]{\linewidth}\raggedright
Variable
\end{minipage} & \begin{minipage}[b]{\linewidth}\raggedright
Interpretation
\end{minipage} \\
\midrule\noalign{}
\endhead
\bottomrule\noalign{}
\endlastfoot
1 & temperature\_c & Moderate variation --- typical indoor growing
temperature, stable across time. \\
2 & humidity\_percent & Reasonably humid environment --- variation
likely due to daily cycles. \\
3 & tds\_value\_ppm & Stable nutrient concentration --- healthy
hydroponic range. \\
4 & p\_h\_level & Excellent --- near-neutral, stable pH. \\
5 & growth\_days & Data spans about 1.5 months: from 2023-08-03 to
2023-09-19. \\
\end{longtable}

\textbf{📊 Insights}

No missing data or extreme anomalies.

Boxplots showed relatively symmetric distributions for most
environmental variables, with slight skewness or mild outliers in TDS
and Growth Days.

Overall, variability within variables supports robust analysis of
growth-environment interactions.

\textbf{🧠 Interpretation}

Each variable behaves consistently with expectations for greenhouse
data:

Environmental factors are well-controlled but still varied enough to
analyze correlations with plant growth.

The dataset is statistically rich and suitable for later modeling (e.g.,
regression or correlation analysis).

\subsection{Stage 3 of EDA: Bivariate/Multivariate analysis
(``Understanding relationships''
phase)}\label{stage-3-of-eda-bivariatemultivariate-analysis-understanding-relationships-phase}

\begin{Shaded}
\begin{Highlighting}[]
\NormalTok{df }\SpecialCharTok{\%\textgreater{}\%} \FunctionTok{select}\NormalTok{(}\FunctionTok{where}\NormalTok{(is.numeric)) }\SpecialCharTok{\%\textgreater{}\%} 
  \FunctionTok{cor}\NormalTok{(}\AttributeTok{use =} \StringTok{"complete.obs"}\NormalTok{) }\SpecialCharTok{\%\textgreater{}\%} \FunctionTok{round}\NormalTok{(}\DecValTok{2}\NormalTok{)}
\end{Highlighting}
\end{Shaded}

\begin{verbatim}
##                  temperature_c humidity_percent tds_value_ppm p_h_level
## temperature_c             1.00             0.03         -0.01      0.01
## humidity_percent          0.03             1.00         -0.01      0.02
## tds_value_ppm            -0.01            -0.01          1.00     -0.01
## p_h_level                 0.01             0.02         -0.01      1.00
## growth_days              -0.07            -0.01         -0.02      0.00
##                  growth_days
## temperature_c          -0.07
## humidity_percent       -0.01
## tds_value_ppm          -0.02
## p_h_level               0.00
## growth_days             1.00
\end{verbatim}

\begin{Shaded}
\begin{Highlighting}[]
\NormalTok{df }\SpecialCharTok{\%\textgreater{}\%} \FunctionTok{select}\NormalTok{(}\FunctionTok{where}\NormalTok{(is.numeric)) }\SpecialCharTok{\%\textgreater{}\%} 
  \FunctionTok{cor}\NormalTok{(}\AttributeTok{use =} \StringTok{"complete.obs"}\NormalTok{) }\SpecialCharTok{\%\textgreater{}\%} 
\NormalTok{  corrplot}\SpecialCharTok{::}\FunctionTok{corrplot}\NormalTok{(}\AttributeTok{method =} \StringTok{"color"}\NormalTok{, }
                     \AttributeTok{type =} \StringTok{"upper"}\NormalTok{, }
                     \AttributeTok{order =} \StringTok{"hclust"}\NormalTok{,}
                     \AttributeTok{addCoef.col =} \StringTok{"black"}\NormalTok{,}
                     \AttributeTok{tl.cex =} \FloatTok{0.8}\NormalTok{,}
                     \AttributeTok{tl.col =} \StringTok{"black"}\NormalTok{,}
                     \AttributeTok{tl.srt =} \DecValTok{45}\NormalTok{,}
                     \AttributeTok{cl.pos =} \StringTok{"n"}\NormalTok{)}
\end{Highlighting}
\end{Shaded}

\pandocbounded{\includegraphics[keepaspectratio]{Lettuce_growth_days_v2.1--Step-2-_files/figure-latex/unnamed-chunk-22-1.pdf}}

\begin{Shaded}
\begin{Highlighting}[]
\FunctionTok{pairs}\NormalTok{(df)  }
\end{Highlighting}
\end{Shaded}

\pandocbounded{\includegraphics[keepaspectratio]{Lettuce_growth_days_v2.1--Step-2-_files/figure-latex/unnamed-chunk-23-1.pdf}}

\begin{Shaded}
\begin{Highlighting}[]
\NormalTok{df }\SpecialCharTok{\%\textgreater{}\%} \FunctionTok{select}\NormalTok{(}\SpecialCharTok{{-}}\NormalTok{plant\_id, }\SpecialCharTok{{-}}\NormalTok{date) }\SpecialCharTok{\%\textgreater{}\%} 
  \FunctionTok{ggpairs}\NormalTok{()}
\end{Highlighting}
\end{Shaded}

\pandocbounded{\includegraphics[keepaspectratio]{Lettuce_growth_days_v2.1--Step-2-_files/figure-latex/unnamed-chunk-24-1.pdf}}
\textbf{🎯 Purpose}

This stage examined how environmental variables relate to each other and
to the target variable --- Growth Days.

🔗 Correlation Matrix (Rounded to 2 decimals)

\begin{longtable}[]{@{}
  >{\raggedright\arraybackslash}p{(\linewidth - 10\tabcolsep) * \real{0.1957}}
  >{\raggedright\arraybackslash}p{(\linewidth - 10\tabcolsep) * \real{0.1739}}
  >{\raggedright\arraybackslash}p{(\linewidth - 10\tabcolsep) * \real{0.2065}}
  >{\raggedright\arraybackslash}p{(\linewidth - 10\tabcolsep) * \real{0.1630}}
  >{\raggedright\arraybackslash}p{(\linewidth - 10\tabcolsep) * \real{0.1196}}
  >{\raggedright\arraybackslash}p{(\linewidth - 10\tabcolsep) * \real{0.1413}}@{}}
\toprule\noalign{}
\begin{minipage}[b]{\linewidth}\raggedright
Variable
\end{minipage} & \begin{minipage}[b]{\linewidth}\raggedright
temperature\_c
\end{minipage} & \begin{minipage}[b]{\linewidth}\raggedright
humidity\_percent
\end{minipage} & \begin{minipage}[b]{\linewidth}\raggedright
tds\_value\_ppm
\end{minipage} & \begin{minipage}[b]{\linewidth}\raggedright
p\_h\_level
\end{minipage} & \begin{minipage}[b]{\linewidth}\raggedright
growth\_days
\end{minipage} \\
\midrule\noalign{}
\endhead
\bottomrule\noalign{}
\endlastfoot
temperature\_c & 1.00 & 0.03 & -0.01 & 0.01 & -0.07 \\
humidity\_percent & 0.03 & 1.00 & -0.01 & 0.02 & -0.01 \\
tds\_value\_ppm & -0.01 & -0.01 & 1.00 & -0.01 & -0.02 \\
p\_h\_level & 0.01 & 0.02 & -0.01 & 1.00 & 0.00 \\
growth\_days & -0.07 & -0.01 & -0.02 & 0.00 & 1.00 \\
\end{longtable}

\textbf{💡 Interpretation of Correlations}

All correlations are weak (\textbar r\textbar{} \textless{} 0.1) --- no
strong linear dependencies among variables.

Slight negative correlation between temperature\_c and growth\_days
(-0.07) could indicate that higher temperatures may slightly accelerate
early growth, reducing total growth duration.

Other environmental factors (humidity\_percent, tds\_value\_ppm,
p\_h\_level) show minimal direct relationships to growth time.

The correlation plot confirms overall independence of predictors,
meaning multicollinearity is not a concern.

\textbf{🧠 General Interpretation}

The weak correlations suggest lettuce growth in this dataset is not
linearly dominated by any single factor --- growth may depend on
nonlinear or combined environmental effects (e.g., optimal ranges rather
than simple linear relationships).

These findings justify future modeling efforts (like regression or
machine learning) to explore interactions or thresholds rather than
direct proportional effects.

\textbf{✅ Summary of the Three Phases}

\begin{longtable}[]{@{}
  >{\raggedright\arraybackslash}p{(\linewidth - 4\tabcolsep) * \real{0.2500}}
  >{\raggedright\arraybackslash}p{(\linewidth - 4\tabcolsep) * \real{0.2500}}
  >{\raggedright\arraybackslash}p{(\linewidth - 4\tabcolsep) * \real{0.5000}}@{}}
\toprule\noalign{}
\begin{minipage}[b]{\linewidth}\raggedright
Phase
\end{minipage} & \begin{minipage}[b]{\linewidth}\raggedright
Focus
\end{minipage} & \begin{minipage}[b]{\linewidth}\raggedright
Key Findings
\end{minipage} \\
\midrule\noalign{}
\endhead
\bottomrule\noalign{}
\endlastfoot
3.1 Data Discovery \& Profiling & Structure \& cleanliness & Data was
complete, consistent, and properly formatted. \\
3.2 Univariate Analysis & Individual variable behavior & Variables
showed natural variation without missing data; distributions were
realistic. \\
3.3 Bivariate/Multivariate Analysis & Variable relationships & No strong
linear relationships; temperature showed a slight negative association
with growth days. \\
\end{longtable}

\begin{center}\rule{0.5\linewidth}{0.5pt}\end{center}

All statistics by plant id

\begin{Shaded}
\begin{Highlighting}[]
\NormalTok{all\_stats\_by\_id }\OtherTok{\textless{}{-}}\NormalTok{ df }\SpecialCharTok{\%\textgreater{}\%} \FunctionTok{group\_by}\NormalTok{(plant\_id) }\SpecialCharTok{\%\textgreater{}\%}
  \FunctionTok{summarise}\NormalTok{(}
            \AttributeTok{date\_min =} \FunctionTok{min}\NormalTok{(date),}
            \AttributeTok{date\_mean =} \FunctionTok{mean}\NormalTok{(date),}
            \AttributeTok{date\_median =} \FunctionTok{median}\NormalTok{(date),}
            \AttributeTok{date\_max =} \FunctionTok{max}\NormalTok{(date),}
            \AttributeTok{temperature\_c\_min =} \FunctionTok{min}\NormalTok{(temperature\_c),}
            \AttributeTok{temperature\_c\_q1 =} \FunctionTok{quantile}\NormalTok{(temperature\_c, }\FloatTok{0.25}\NormalTok{),}
            \AttributeTok{temperature\_c\_mean =} \FunctionTok{mean}\NormalTok{(temperature\_c),}
            \AttributeTok{temperature\_c\_median =} \FunctionTok{median}\NormalTok{(temperature\_c),}
            \AttributeTok{temperature\_c\_sd =} \FunctionTok{sd}\NormalTok{(temperature\_c),}
            \AttributeTok{temperature\_c\_q3 =} \FunctionTok{quantile}\NormalTok{(temperature\_c, }\FloatTok{0.75}\NormalTok{),}
            \AttributeTok{temperature\_c\_max =} \FunctionTok{max}\NormalTok{(temperature\_c),}
            \AttributeTok{humidity\_percent\_min =} \FunctionTok{min}\NormalTok{(humidity\_percent),}
            \AttributeTok{humidity\_percent\_q1 =} \FunctionTok{quantile}\NormalTok{(humidity\_percent, }\FloatTok{0.25}\NormalTok{),}
            \AttributeTok{humidity\_percent\_mean =} \FunctionTok{mean}\NormalTok{(humidity\_percent),}
            \AttributeTok{humidity\_percent\_median =} \FunctionTok{median}\NormalTok{(humidity\_percent),}
            \AttributeTok{humidity\_percent\_sd =} \FunctionTok{sd}\NormalTok{(humidity\_percent),}
            \AttributeTok{humidity\_percent\_q3 =} \FunctionTok{quantile}\NormalTok{(humidity\_percent, }\FloatTok{0.75}\NormalTok{),}
            \AttributeTok{humidity\_percent\_max =} \FunctionTok{max}\NormalTok{(humidity\_percent),}
            \AttributeTok{tds\_value\_ppm\_min =} \FunctionTok{min}\NormalTok{(tds\_value\_ppm),}
            \AttributeTok{tds\_value\_ppm\_q1 =} \FunctionTok{quantile}\NormalTok{(tds\_value\_ppm, }\FloatTok{0.25}\NormalTok{),}
            \AttributeTok{tds\_value\_ppm\_mean =} \FunctionTok{mean}\NormalTok{(tds\_value\_ppm),}
            \AttributeTok{tds\_value\_ppm\_median =} \FunctionTok{median}\NormalTok{(tds\_value\_ppm),}
            \AttributeTok{tds\_value\_ppm\_sd =} \FunctionTok{sd}\NormalTok{(tds\_value\_ppm),}
            \AttributeTok{tds\_value\_ppm\_q3 =} \FunctionTok{quantile}\NormalTok{(tds\_value\_ppm, }\FloatTok{0.75}\NormalTok{),}
            \AttributeTok{tds\_value\_ppm\_max =} \FunctionTok{max}\NormalTok{(tds\_value\_ppm),}
            \AttributeTok{p\_h\_level\_min =} \FunctionTok{min}\NormalTok{(p\_h\_level),}
            \AttributeTok{p\_h\_level\_q1 =} \FunctionTok{quantile}\NormalTok{(p\_h\_level, }\FloatTok{0.25}\NormalTok{),}
            \AttributeTok{p\_h\_level\_mean =} \FunctionTok{mean}\NormalTok{(p\_h\_level),}
            \AttributeTok{p\_h\_level\_median =} \FunctionTok{median}\NormalTok{(p\_h\_level),}
            \AttributeTok{p\_h\_level\_sd =} \FunctionTok{sd}\NormalTok{(p\_h\_level),}
            \AttributeTok{p\_h\_level\_q3 =} \FunctionTok{quantile}\NormalTok{(p\_h\_level, }\FloatTok{0.75}\NormalTok{),}
            \AttributeTok{p\_h\_level\_max =} \FunctionTok{max}\NormalTok{(p\_h\_level),}
            \AttributeTok{growth\_days\_min =} \FunctionTok{min}\NormalTok{(growth\_days),}
            \AttributeTok{growth\_days\_q1 =} \FunctionTok{quantile}\NormalTok{(growth\_days, }\FloatTok{0.25}\NormalTok{),}
            \AttributeTok{growth\_days\_mean =} \FunctionTok{mean}\NormalTok{(growth\_days),}
            \AttributeTok{growth\_days\_median =} \FunctionTok{median}\NormalTok{(growth\_days),}
            \AttributeTok{growth\_days\_sd =} \FunctionTok{sd}\NormalTok{(growth\_days),}
            \AttributeTok{growth\_days\_q3 =} \FunctionTok{quantile}\NormalTok{(growth\_days, }\FloatTok{0.75}\NormalTok{),}
            \AttributeTok{growth\_days\_max =} \FunctionTok{max}\NormalTok{(growth\_days)}
\NormalTok{  ) }\SpecialCharTok{\%\textgreater{}\%} 
  \FunctionTok{mutate}\NormalTok{(}
    \FunctionTok{across}\NormalTok{(}\FunctionTok{where}\NormalTok{(is.numeric), }\SpecialCharTok{\textasciitilde{}} \FunctionTok{round}\NormalTok{(.x, }\DecValTok{2}\NormalTok{))}
\NormalTok{         )}

\NormalTok{all\_stats\_by\_id}
\end{Highlighting}
\end{Shaded}

\begin{verbatim}
## # A tibble: 70 x 40
##    plant_id date_min   date_mean  date_median date_max   temperature_c_min
##    <fct>    <date>     <date>     <date>      <date>                 <dbl>
##  1 1        2023-08-03 2023-08-25 2023-08-25  2023-09-16              20.1
##  2 2        2023-08-03 2023-08-25 2023-08-25  2023-09-16              20.1
##  3 3        2023-08-03 2023-08-26 2023-08-26  2023-09-18              20.1
##  4 4        2023-08-03 2023-08-26 2023-08-26  2023-09-19              20.1
##  5 5        2023-08-03 2023-08-25 2023-08-25  2023-09-16              20.1
##  6 6        2023-08-03 2023-08-25 2023-08-25  2023-09-16              20.1
##  7 7        2023-08-03 2023-08-26 2023-08-26  2023-09-18              20.1
##  8 8        2023-08-03 2023-08-25 2023-08-25  2023-09-16              20.1
##  9 9        2023-08-03 2023-08-25 2023-08-25  2023-09-17              20.1
## 10 10       2023-08-03 2023-08-25 2023-08-25  2023-09-16              20.1
## # i 60 more rows
## # i 34 more variables: temperature_c_q1 <dbl>, temperature_c_mean <dbl>,
## #   temperature_c_median <dbl>, temperature_c_sd <dbl>, temperature_c_q3 <dbl>,
## #   temperature_c_max <dbl>, humidity_percent_min <dbl>,
## #   humidity_percent_q1 <dbl>, humidity_percent_mean <dbl>,
## #   humidity_percent_median <dbl>, humidity_percent_sd <dbl>,
## #   humidity_percent_q3 <dbl>, humidity_percent_max <dbl>, ...
\end{verbatim}

Summary of all statistics by plant id

\begin{Shaded}
\begin{Highlighting}[]
\NormalTok{all\_stats\_by\_id }\SpecialCharTok{\%\textgreater{}\%} \FunctionTok{summary}\NormalTok{()}
\end{Highlighting}
\end{Shaded}

\begin{verbatim}
##     plant_id     date_min            date_mean           date_median        
##  1      : 1   Min.   :2023-08-03   Min.   :2023-08-25   Min.   :2023-08-25  
##  2      : 1   1st Qu.:2023-08-03   1st Qu.:2023-08-25   1st Qu.:2023-08-25  
##  3      : 1   Median :2023-08-03   Median :2023-08-25   Median :2023-08-25  
##  4      : 1   Mean   :2023-08-03   Mean   :2023-08-25   Mean   :2023-08-25  
##  5      : 1   3rd Qu.:2023-08-03   3rd Qu.:2023-08-25   3rd Qu.:2023-08-25  
##  6      : 1   Max.   :2023-08-03   Max.   :2023-08-26   Max.   :2023-08-26  
##  (Other):64                                                                 
##     date_max          temperature_c_min temperature_c_q1 temperature_c_mean
##  Min.   :2023-09-16   Min.   :18.00     Min.   :18.80    Min.   :21.04     
##  1st Qu.:2023-09-16   1st Qu.:18.32     1st Qu.:20.25    1st Qu.:21.94     
##  Median :2023-09-16   Median :20.10     Median :30.10    Median :30.63     
##  Mean   :2023-09-16   Mean   :19.56     Mean   :27.27    Mean   :28.14     
##  3rd Qu.:2023-09-16   3rd Qu.:20.10     3rd Qu.:30.10    3rd Qu.:30.63     
##  Max.   :2023-09-19   Max.   :20.10     Max.   :30.10    Max.   :30.67     
##                                                                            
##  temperature_c_median temperature_c_sd temperature_c_q3 temperature_c_max
##  Min.   :20.75        Min.   :1.810    Min.   :22.2     Min.   :24.60    
##  1st Qu.:22.23        1st Qu.:2.277    1st Qu.:23.8     1st Qu.:26.65    
##  Median :30.90        Median :2.380    Median :31.7     Median :33.50    
##  Mean   :28.34        Mean   :2.301    Mean   :29.4     Mean   :31.25    
##  3rd Qu.:30.90        3rd Qu.:2.380    3rd Qu.:31.7     3rd Qu.:33.50    
##  Max.   :31.10        Max.   :3.160    Max.   :31.7     Max.   :33.50    
##                                                                          
##  humidity_percent_min humidity_percent_q1 humidity_percent_mean
##  Min.   :50.00        Min.   :53.00       Min.   :62.62        
##  1st Qu.:50.00        1st Qu.:56.00       1st Qu.:63.71        
##  Median :50.00        Median :57.00       Median :64.82        
##  Mean   :50.19        Mean   :57.39       Mean   :64.88        
##  3rd Qu.:50.00        3rd Qu.:59.00       3rd Qu.:65.81        
##  Max.   :51.00        Max.   :62.50       Max.   :68.68        
##                                                                
##  humidity_percent_median humidity_percent_sd humidity_percent_q3
##  Min.   :61.00           Min.   : 7.050      Min.   :68.00      
##  1st Qu.:63.00           1st Qu.: 8.533      1st Qu.:71.00      
##  Median :65.00           Median : 9.060      Median :72.88      
##  Mean   :64.89           Mean   : 8.959      Mean   :72.29      
##  3rd Qu.:66.00           3rd Qu.: 9.287      3rd Qu.:74.00      
##  Max.   :72.00           Max.   :10.290      Max.   :76.00      
##                                                                 
##  humidity_percent_max tds_value_ppm_min tds_value_ppm_q1 tds_value_ppm_mean
##  Min.   :77.00        Min.   :400       Min.   :451.0    Min.   :567.0     
##  1st Qu.:79.00        1st Qu.:402       1st Qu.:486.2    1st Qu.:584.2     
##  Median :80.00        Median :405       Median :501.0    Median :598.3     
##  Mean   :79.61        Mean   :409       Mean   :502.7    Mean   :598.0     
##  3rd Qu.:80.00        3rd Qu.:412       3rd Qu.:515.8    3rd Qu.:611.9     
##  Max.   :80.00        Max.   :449       Max.   :576.0    Max.   :628.2     
##                                                                            
##  tds_value_ppm_median tds_value_ppm_sd tds_value_ppm_q3 tds_value_ppm_max
##  Min.   :544.0        Min.   : 95.55   Min.   :641.0    Min.   :765.0    
##  1st Qu.:574.2        1st Qu.:109.97   1st Qu.:680.0    1st Qu.:789.0    
##  Median :596.5        Median :115.61   Median :698.8    Median :795.0    
##  Mean   :596.9        Mean   :115.56   Mean   :695.0    Mean   :791.9    
##  3rd Qu.:613.0        3rd Qu.:120.48   3rd Qu.:714.2    3rd Qu.:798.0    
##  Max.   :666.0        Max.   :138.25   Max.   :749.0    Max.   :800.0    
##                                                                          
##  p_h_level_min    p_h_level_q1   p_h_level_mean  p_h_level_median
##  Min.   :6.000   Min.   :6.100   Min.   :6.310   Min.   :6.300   
##  1st Qu.:6.000   1st Qu.:6.200   1st Qu.:6.380   1st Qu.:6.400   
##  Median :6.000   Median :6.200   Median :6.400   Median :6.400   
##  Mean   :6.001   Mean   :6.207   Mean   :6.399   Mean   :6.398   
##  3rd Qu.:6.000   3rd Qu.:6.200   3rd Qu.:6.420   3rd Qu.:6.400   
##  Max.   :6.100   Max.   :6.400   Max.   :6.490   Max.   :6.500   
##                                                                  
##   p_h_level_sd     p_h_level_q3   p_h_level_max   growth_days_min
##  Min.   :0.1900   Min.   :6.500   Min.   :6.700   Min.   :1      
##  1st Qu.:0.2200   1st Qu.:6.600   1st Qu.:6.800   1st Qu.:1      
##  Median :0.2300   Median :6.600   Median :6.800   Median :1      
##  Mean   :0.2344   Mean   :6.598   Mean   :6.789   Mean   :1      
##  3rd Qu.:0.2500   3rd Qu.:6.600   3rd Qu.:6.800   3rd Qu.:1      
##  Max.   :0.2800   Max.   :6.700   Max.   :6.800   Max.   :1      
##                                                                  
##  growth_days_q1  growth_days_mean growth_days_median growth_days_sd 
##  Min.   :12.00   Min.   :23.00    Min.   :23.00      Min.   :13.13  
##  1st Qu.:12.00   1st Qu.:23.00    1st Qu.:23.00      1st Qu.:13.13  
##  Median :12.00   Median :23.00    Median :23.00      Median :13.13  
##  Mean   :12.07   Mean   :23.13    Mean   :23.14      Mean   :13.21  
##  3rd Qu.:12.00   3rd Qu.:23.00    3rd Qu.:23.00      3rd Qu.:13.13  
##  Max.   :12.75   Max.   :24.50    Max.   :24.50      Max.   :14.00  
##                                                                     
##  growth_days_q3  growth_days_max
##  Min.   :34.00   Min.   :45.00  
##  1st Qu.:34.00   1st Qu.:45.00  
##  Median :34.00   Median :45.00  
##  Mean   :34.20   Mean   :45.26  
##  3rd Qu.:34.00   3rd Qu.:45.00  
##  Max.   :36.25   Max.   :48.00  
## 
\end{verbatim}

Statistics (basic) by plant id

\begin{Shaded}
\begin{Highlighting}[]
\NormalTok{all\_stats\_by\_id }\SpecialCharTok{\%\textgreater{}\%} 
  \FunctionTok{select}\NormalTok{(plant\_id, }
\NormalTok{         temperature\_c\_mean,}
\NormalTok{         humidity\_percent\_mean,}
\NormalTok{         tds\_value\_ppm\_mean,}
\NormalTok{         p\_h\_level\_mean,}
\NormalTok{         growth\_days\_max)}
\end{Highlighting}
\end{Shaded}

\begin{verbatim}
## # A tibble: 70 x 6
##    plant_id temperature_c_mean humidity_percent_mean tds_value_ppm_mean
##    <fct>                 <dbl>                 <dbl>              <dbl>
##  1 1                      30.6                  64.0               615.
##  2 2                      30.6                  64.6               597.
##  3 3                      30.6                  68.7               620.
##  4 4                      30.6                  63.4               598.
##  5 5                      30.6                  65.1               577.
##  6 6                      30.6                  64.4               614.
##  7 7                      30.6                  64.1               608.
##  8 8                      30.6                  66.3               598.
##  9 9                      30.6                  66.9               608.
## 10 10                     30.6                  65.8               613.
## # i 60 more rows
## # i 2 more variables: p_h_level_mean <dbl>, growth_days_max <dbl>
\end{verbatim}

Statistics by maximum growth day

\begin{Shaded}
\begin{Highlighting}[]
\NormalTok{df }\SpecialCharTok{\%\textgreater{}\%} \FunctionTok{group\_by}\NormalTok{(plant\_id) }\SpecialCharTok{\%\textgreater{}\%} 
  \FunctionTok{mutate}\NormalTok{(}\AttributeTok{growth\_days\_max =} \FunctionTok{max}\NormalTok{(growth\_days)) }\SpecialCharTok{\%\textgreater{}\%} 
  \FunctionTok{group\_by}\NormalTok{(growth\_days\_max) }\SpecialCharTok{\%\textgreater{}\%} 
  \FunctionTok{summarise}\NormalTok{(}
    \AttributeTok{plant\_id\_n =} \FunctionTok{n\_distinct}\NormalTok{(plant\_id),}
            \AttributeTok{temperature\_c\_mean =} \FunctionTok{mean}\NormalTok{(temperature\_c),}
            \AttributeTok{humidity\_percent\_mean =} \FunctionTok{mean}\NormalTok{(humidity\_percent),}
            \AttributeTok{tds\_value\_ppm\_mean =} \FunctionTok{mean}\NormalTok{(tds\_value\_ppm),}
            \AttributeTok{p\_h\_level\_mean =} \FunctionTok{mean}\NormalTok{(p\_h\_level)}
\NormalTok{  ) }\SpecialCharTok{\%\textgreater{}\%} 
  \FunctionTok{mutate}\NormalTok{(}
    \FunctionTok{across}\NormalTok{(}
      \FunctionTok{where}\NormalTok{(is.numeric), }\SpecialCharTok{\textasciitilde{}} \FunctionTok{round}\NormalTok{(.x, }\DecValTok{2}\NormalTok{)}
\NormalTok{      )}
\NormalTok{  )}
\end{Highlighting}
\end{Shaded}

\begin{verbatim}
## # A tibble: 4 x 6
##   growth_days_max plant_id_n temperature_c_mean humidity_percent_mean
##             <dbl>      <dbl>              <dbl>                 <dbl>
## 1              45         59               28.1                  64.9
## 2              46          5               27.0                  65.0
## 3              47          5               28.8                  65.0
## 4              48          1               30.6                  63.4
## # i 2 more variables: tds_value_ppm_mean <dbl>, p_h_level_mean <dbl>
\end{verbatim}

\section{3.4 Visualize relationships}\label{visualize-relationships}

Time-series plots (trends over time)

\begin{Shaded}
\begin{Highlighting}[]
\NormalTok{df }\SpecialCharTok{\%\textgreater{}\%}
  \FunctionTok{pivot\_longer}\NormalTok{(}\SpecialCharTok{{-}}\FunctionTok{c}\NormalTok{(date, plant\_id), }\AttributeTok{names\_to =} \StringTok{"variable"}\NormalTok{, }\AttributeTok{values\_to =} \StringTok{"value"}\NormalTok{) }\SpecialCharTok{\%\textgreater{}\%} 
  \FunctionTok{group\_by}\NormalTok{(variable, date) }\SpecialCharTok{\%\textgreater{}\%}
  \FunctionTok{summarise}\NormalTok{(}
    \AttributeTok{mean =} \FunctionTok{mean}\NormalTok{(value, }\AttributeTok{na.rm =} \ConstantTok{TRUE}\NormalTok{),}
    \AttributeTok{p25  =} \FunctionTok{quantile}\NormalTok{(value, }\FloatTok{0.25}\NormalTok{, }\AttributeTok{na.rm =} \ConstantTok{TRUE}\NormalTok{),}
    \AttributeTok{p75  =} \FunctionTok{quantile}\NormalTok{(value, }\FloatTok{0.75}\NormalTok{, }\AttributeTok{na.rm =} \ConstantTok{TRUE}\NormalTok{),}
    \AttributeTok{.groups =} \StringTok{"drop"}
\NormalTok{  ) }\SpecialCharTok{\%\textgreater{}\%}
  \FunctionTok{arrange}\NormalTok{(variable, date) }\SpecialCharTok{\%\textgreater{}\%} 
  \FunctionTok{ggplot}\NormalTok{(}\FunctionTok{aes}\NormalTok{(date, mean)) }\SpecialCharTok{+}
  \FunctionTok{geom\_ribbon}\NormalTok{(}\FunctionTok{aes}\NormalTok{(}\AttributeTok{ymin =}\NormalTok{ p25, }\AttributeTok{ymax =}\NormalTok{ p75), }\AttributeTok{alpha =} \FloatTok{0.15}\NormalTok{) }\SpecialCharTok{+}
  \FunctionTok{geom\_line}\NormalTok{() }\SpecialCharTok{+}
  \FunctionTok{facet\_wrap}\NormalTok{(}\SpecialCharTok{\textasciitilde{}}\NormalTok{ variable, }\AttributeTok{scales =} \StringTok{"free\_y"}\NormalTok{, }\AttributeTok{ncol =} \DecValTok{2}\NormalTok{) }\SpecialCharTok{+}
  \FunctionTok{scale\_x\_date}\NormalTok{(}\AttributeTok{date\_breaks =} \StringTok{"1 month"}\NormalTok{, }\AttributeTok{date\_labels =} \StringTok{"\%b \%Y"}\NormalTok{) }\SpecialCharTok{+}
  \FunctionTok{labs}\NormalTok{(}\AttributeTok{x =} \ConstantTok{NULL}\NormalTok{, }\AttributeTok{y =} \ConstantTok{NULL}\NormalTok{, }\AttributeTok{title =} \StringTok{"Daily mean with IQR band"}\NormalTok{) }\SpecialCharTok{+}
  \FunctionTok{theme\_minimal}\NormalTok{() }\SpecialCharTok{+}
  \FunctionTok{theme}\NormalTok{(}\AttributeTok{axis.text.x =} \FunctionTok{element\_text}\NormalTok{(}\AttributeTok{angle =} \DecValTok{45}\NormalTok{, }\AttributeTok{hjust =} \DecValTok{1}\NormalTok{),}
        \AttributeTok{panel.grid.minor =} \FunctionTok{element\_blank}\NormalTok{())}
\end{Highlighting}
\end{Shaded}

\pandocbounded{\includegraphics[keepaspectratio]{Lettuce_growth_days_v2.1--Step-2-_files/figure-latex/unnamed-chunk-29-1.pdf}}
c

\begin{Shaded}
\begin{Highlighting}[]
\NormalTok{df }\SpecialCharTok{\%\textgreater{}\%} 
  \FunctionTok{pivot\_longer}\NormalTok{(}\SpecialCharTok{{-}}\FunctionTok{c}\NormalTok{(date, plant\_id), }\AttributeTok{names\_to =} \StringTok{"variable"}\NormalTok{, }\AttributeTok{values\_to =} \StringTok{"values"}\NormalTok{) }\SpecialCharTok{\%\textgreater{}\%} 
  \FunctionTok{group\_by}\NormalTok{(variable, date) }\SpecialCharTok{\%\textgreater{}\%} 
  \FunctionTok{summarise}\NormalTok{(}\AttributeTok{mean =} \FunctionTok{mean}\NormalTok{(values)) }\SpecialCharTok{\%\textgreater{}\%} 
  \FunctionTok{arrange}\NormalTok{(variable, date) }\SpecialCharTok{\%\textgreater{}\%} 
  \FunctionTok{ggplot}\NormalTok{(}\FunctionTok{aes}\NormalTok{(}\AttributeTok{x =}\NormalTok{ date, }\AttributeTok{y =}\NormalTok{ mean)) }\SpecialCharTok{+}
  \FunctionTok{geom\_line}\NormalTok{() }\SpecialCharTok{+} 
  \FunctionTok{facet\_wrap}\NormalTok{(}\SpecialCharTok{\textasciitilde{}}\NormalTok{ variable, }\AttributeTok{scales =} \StringTok{"free\_y"}\NormalTok{, }\AttributeTok{ncol =} \DecValTok{2}\NormalTok{) }\SpecialCharTok{+} 
  \FunctionTok{scale\_x\_date}\NormalTok{(}\AttributeTok{date\_breaks =} \StringTok{"1 month"}\NormalTok{, }\AttributeTok{date\_labels =} \StringTok{"\%b \%Y"}\NormalTok{) }\SpecialCharTok{+}
  \FunctionTok{labs}\NormalTok{(}\AttributeTok{x =} \ConstantTok{NULL}\NormalTok{, }\AttributeTok{y =} \ConstantTok{NULL}\NormalTok{, }\AttributeTok{title =} \StringTok{"Daily mean"}\NormalTok{) }\SpecialCharTok{+} 
  \CommentTok{\# theme\_minimal() +}
  \FunctionTok{theme}\NormalTok{(}\AttributeTok{axis.text.x =} \FunctionTok{element\_text}\NormalTok{(}\AttributeTok{angle =} \DecValTok{45}\NormalTok{, }\AttributeTok{hjust =} \DecValTok{1}\NormalTok{),}
        \AttributeTok{panel.grid.minor =} \FunctionTok{element\_blank}\NormalTok{())}
\end{Highlighting}
\end{Shaded}

\begin{verbatim}
## `summarise()` has regrouped the output.
## i Summaries were computed grouped by variable and date.
## i Output is grouped by variable.
## i Use `summarise(.groups = "drop_last")` to silence this message.
## i Use `summarise(.by = c(variable, date))` for per-operation grouping
##   (`?dplyr::dplyr_by`) instead.
\end{verbatim}

\pandocbounded{\includegraphics[keepaspectratio]{Lettuce_growth_days_v2.1--Step-2-_files/figure-latex/unnamed-chunk-30-1.pdf}}
Scatterplots (check correlation)

\begin{Shaded}
\begin{Highlighting}[]
\NormalTok{df }\SpecialCharTok{\%\textgreater{}\%} \FunctionTok{select}\NormalTok{(}\FunctionTok{where}\NormalTok{(is.numeric)) }\SpecialCharTok{\%\textgreater{}\%} 
  \FunctionTok{ggpairs}\NormalTok{(}\AttributeTok{progress =} \ConstantTok{FALSE}\NormalTok{,}
          \AttributeTok{upper =} \FunctionTok{list}\NormalTok{(}\AttributeTok{continuous =} \FunctionTok{wrap}\NormalTok{(}\StringTok{"cor"}\NormalTok{, }\AttributeTok{method =} \StringTok{"pearson"}\NormalTok{)),}
          \AttributeTok{diag =} \FunctionTok{list}\NormalTok{(}\AttributeTok{continuous =} \StringTok{"densityDiag"}\NormalTok{),}
          \AttributeTok{lwer =} \FunctionTok{list}\NormalTok{(}\AttributeTok{continuous =} \FunctionTok{wrap}\NormalTok{(}\StringTok{"points"}\NormalTok{, }\AttributeTok{alpha =} \FloatTok{0.4}\NormalTok{, }\AttributeTok{size =} \FloatTok{0.7}\NormalTok{))}
    
\NormalTok{  )}
\end{Highlighting}
\end{Shaded}

\begin{verbatim}
## Warning: The `...` argument of `ggpais()` is deprecated as of GGally 2.3.0.
## This warning is displayed once per session.
## Call `lifecycle::last_lifecycle_warnings()` to see where this warning was
## generated.
\end{verbatim}

\pandocbounded{\includegraphics[keepaspectratio]{Lettuce_growth_days_v2.1--Step-2-_files/figure-latex/unnamed-chunk-31-1.pdf}}

\section{3.5 Form hypotheses}\label{form-hypotheses}

Does higher humidity improve growth?

Are weekly growth patterns consistent?

Is there seasonality in the data?

\end{document}
